\subsection{向量空间上的广义函数与缓增广义函数}

设 \(u\) 是从 \(\mathbf{R}^{n}\) 到 \(\mathbf{R}^{n}\) 中的双射线性映射。映射 \(\varphi \mapsto \varphi \circ u\) 是 \(\mathcal{S}(\mathbf{R}^{n})\) （相应地， \(\mathcal{D}(\mathbf{R}^{n})\) ）的自同构；它的转置是 \(\mathcal{S}'(\mathbf{R}^{n})\) （相应地， \(\mathcal{D}'(\mathbf{R}^{n})\) ）的自同构。

设 E 是 n 维实向量空间。设 \(v: \mathbf{R}^{n} \to E\) 是向量空间同构。我们用 \(\mathcal{S}(E)\) （相应地， \(\mathcal{D}(E)\) ）表示映射 \(\varphi: E \to \mathbf{C}\) 中满足 \(\varphi \circ v \in \mathcal{S}(\mathbf{R}^{n})\) （相应地，满足 \(\varphi \circ v \in \mathcal{D}(\mathbf{R}^{n})\) ）者所成的集。由前面的注，这些空间不依赖于 v 的选取；它们同构于 \(\mathcal{S}(\mathbf{R}^{n})\) （相应地， \(\mathcal{D}(\mathbf{R}^{n})\) ）。我们用 \(\mathcal{S}'(E)\) （相应地， \(\mathcal{D}'(E)\) ）表示 \(\mathcal{S}(E)\) （相应地， \(\mathcal{D}(E)\) ）的对偶，赋有有界收敛拓扑。它是同构于 \(\mathcal{S}'(\mathbf{R}^{n})\) （相应地，同构于 \(\mathcal{D}'(\mathbf{R}^{n})\) ）的拓扑向量空间。

设 E 与 F 是 n 维实向量空间，它们关于双线性形式 \(b: E \times F \to \mathbf{R}\) 处于对偶。局部紧交换群 E 关于映射 

\[
(x, y) \mapsto \exp (2 i \pi b (x, y))
\]

（从 E × F 到 U 中；参见 II, p. 235 的推论 1）与 F 处于对偶。我们给 E 与 F 赋以关于这一映射互为对偶的哈尔测度。

空间 \(\mathcal{S}(\mathrm{E})\) 包含在 \(\mathrm{L}^1 (\mathrm{E})\) 中； E 的傅里叶变换通过过渡到子空间与对偶，诱导出从 \(\mathcal{S}(\mathrm{E})\) 到 \(\mathcal{S}(\mathrm{F})\) 中的拓扑向量空间同构，其转置是从 \(\mathcal{S}'(\mathrm{F})\) 到 \(\mathcal{S}'(\mathrm{E})\) 中的拓扑向量空间同构。

